% CLAUDE REPORT::WR-BRIDGE-0001::2026-11-13T18:58-08:00
% Scheduled user routine "worst-case derivation complexity W_R(N)" (five-assignment shared objective).
% Additive report. Changes no schedules, task count, role ownership or root main.tex region.
\documentclass[11pt]{article}
\usepackage{amsmath,amssymb,amsthm}
\usepackage[margin=1in]{geometry}
\usepackage{booktabs}
\usepackage{hyperref}
\newtheorem{theorem}{Theorem}
\newtheorem{lemma}[theorem]{Lemma}
\newtheorem{corollary}[theorem]{Corollary}
\theoremstyle{definition}
\newtheorem{definition}[theorem]{Definition}
\newtheorem{remark}[theorem]{Remark}
\newcommand{\SAT}{\mathrm{SAT}}
\newcommand{\UNSAT}{\mathrm{UNSAT}}
\newcommand{\Wb}{\overline W}
\title{The Oracle-Collapse Barrier for the Worst-Case Derivation Bridge\\
\large $\SAT\in\mathsf P\Rightarrow \Wb_R(N)\le\mathrm{poly}(N)$ must open the solver's code}
\author{Claude scheduled routine, WR-BRIDGE-0001 (2026-11-13 18:58 PST)}
\date{}
\begin{document}
\maketitle

\noindent\textbf{Status summary.}
Theorem~\ref{thm:collapse} (oracle collapse) and Corollary~\ref{cor:dichotomy}
are \textbf{PROVED/EXACT} (short, elementary; independent review pending).
Lemma~\ref{lem:free} is \textbf{PROVED/EXACT} and records the user's
free-optimal-sequence point.
The exponential resolution lower bound for the resolution-type
multisign hypergraph calculus is a \textbf{KNOWN THEOREM}.
The EF white-box route is \textbf{CONDITIONAL} and already in the corpus (LSPV-1/357, /361).
The Tseitin width table is \textbf{COMPUTATIONAL}.
The bridge for every fixed $R$, a superpolynomial lower bound for any $R$ strong enough
to simulate Gaussian elimination, and \textbf{P versus NP remain OPEN}.

\section{Objective and the free-sequence convention}
For a Cook--Reckhow refutation system $R$ (polynomial-time verifier, sound, complete) with
native cost $\operatorname{cost}_R$ satisfying $\operatorname{cost}_R(\pi)\ge|\pi|^{\Omega(1)}$, put
\[
 w_R(F)=\min_{\pi:F\vdash_R\bot}\operatorname{cost}_R(\pi),\qquad
 \Wb_R(N)=\max_{F\in\UNSAT,\ |\langle F\rangle|\le N} w_R(F).
\]
\[
 \mathrm{Bridge}(R):\qquad \SAT\in\mathsf P\ \Longrightarrow\ \exists k\ \forall N:\ \Wb_R(N)\le (N+2)^k .
\]

\begin{lemma}[free optimal sequence is a valid lower bound]\label{lem:free}
Let $M$ be any procedure that, on every $F\in\UNSAT$, writes an $R$-refutation $\pi_F$ of $F$
on an output tape. Then $T_M(F)\ge|\pi_F|$ and hence
$\max_{|F|\le N}T_M(F)\ge \max_{|F|\le N,F\in\UNSAT}\min_\pi|\pi|$.
Granting $M$ free knowledge of the optimal sequence can only decrease its cost, so the minimum is a
sound lower bound on every derivation-producing procedure. \emph{Proof:} each output symbol costs a
transition; minimizing over $\pi$ is below the particular $\pi_F$. \qed
\end{lemma}

\begin{remark}[what the free choice turns $\Wb_R$ into]
Because the best derivation is granted for free, $\Wb_R$ is a \emph{nondeterministic}
(proof-complexity) measure. Cook--Reckhow~\cite{CR79}: some $R$ has $\Wb_R$ polynomial iff
$\mathsf{NP}=\mathsf{coNP}$. Hence no lower bound on $\Wb_R$ alone can say anything about
\emph{deterministic} time beyond $\mathsf{NP}\neq\mathsf{coNP}$-type content; all P-versus-NP
content lives in $\mathrm{Bridge}(R)$. This is consistent with Lemma~\ref{lem:free}: a SAT
decider is not obliged to output $R$-derivations.
\end{remark}

\section{Main new result: oracle collapse}
For an oracle $A$, define the relativized hypothesis $\SAT\in\mathsf P^A$
(ordinary, unrelativized SAT decided by a polynomial-time oracle machine) and
\[
 \mathrm{Bridge}^A(R):\quad \SAT\in\mathsf P^A\Longrightarrow \Wb_R(N)\le\mathrm{poly}(N).
\]
The conclusion does not mention $A$: the verifier of $R$ is oracle-free, so $\Wb_R$ is the same
absolute integer sequence for every $A$.

\begin{theorem}[oracle collapse]\label{thm:collapse}
For every fixed oracle-free Cook--Reckhow system $R$ the following are equivalent:
(i) $\mathrm{Bridge}^A(R)$ holds for every oracle $A$;
(ii) $\mathrm{Bridge}^{\SAT}(R)$ holds;
(iii) $\Wb_R$ is polynomially bounded ($R$ is p-bounded).
In particular (i) implies $\mathsf{NP}=\mathsf{coNP}$.
\end{theorem}
\begin{proof}
(i)$\Rightarrow$(ii) is specialization. (ii)$\Rightarrow$(iii): $\SAT\in\mathsf P^{\SAT}$ holds
trivially (one query, linear time), so the premise of $\mathrm{Bridge}^{\SAT}(R)$ is true and its
conclusion follows. (iii)$\Rightarrow$(i): the conclusion holds outright.
Finally a p-bounded $R$ puts $\UNSAT\in\mathsf{NP}$, so $\mathsf{NP}=\mathsf{coNP}$~\cite{CR79}. The
same argument works with $A=\mathrm{TQBF}$, because $\mathsf{NP}\subseteq\mathsf{PSPACE}=\mathsf P^{\mathrm{TQBF}}$.
\end{proof}

\begin{corollary}[black-box bridges are self-defeating]\label{cor:dichotomy}
Call a proof of $\mathrm{Bridge}(R)$ \emph{black-box in the decider} if it uses the hypothetical
decider $D$ only through (a) correctness of its input--output function and (b) a polynomial bound on
its step count in a cost model where an oracle query is one step. Every such proof also proves
$\mathrm{Bridge}^{\SAT}(R)$, hence proves $R$ p-bounded and $\mathsf{NP}=\mathsf{coNP}$.
Therefore the program
\[
 \text{``}\mathrm{Bridge}(R)\ \wedge\ \Wb_R\ \text{superpolynomial''}\ \Longrightarrow\ \mathsf P\ne\mathsf{NP}
\]
can succeed only with a proof of $\mathrm{Bridge}(R)$ that (1) does not relativize, and (2) uses
the \emph{code} of the polynomial-time decider (white-box), translating its actual transitions into
$R$-inferences. Free knowledge of the answer, which is harmless on the derivation side (Lemma~\ref{lem:free}),
is fatal on the solver side.
\end{corollary}

\begin{remark}[scope]
This is a statement about proofs of the bridge, not a claim that the bridge is false or unprovable.
It refines LSPV-1/357's informal label ``nonrelativizing sufficient criterion'' into an exact
equivalence. Algebrization~\cite{AW09}: an algebraic extension $\widetilde{\mathrm{TQBF}}$ agrees with
TQBF on Boolean points, so $\SAT\in\mathsf P^{\widetilde{\mathrm{TQBF}}}$ as well and the same collapse
applies to any proof that remains valid for oracle machines with access to such an extension
(\textbf{CONDITIONAL} on adopting the Aaronson--Wigderson definition for one-sided implications).
\end{remark}

\section{Consequences for the five assignments}
\subsection*{Multisign DAG hypergraph calculus (Builder, Hypergraph roles)}
If the reduction rule merges two signed hyperedges that clash in exactly one vertex, the calculus is DAG
resolution. For the odd-charge Tseitin 3-CNF on a 3-regular expander, resolution needs
$2^{\Omega(N)}$ size~\cite{Urq87,BSW01} (\textbf{KNOWN THEOREM}). So $\Wb_{\mathrm{Res}}$ is already
exponential and, by the observation already recorded in the Codex 11-11 audit, $\mathrm{Bridge}(\mathrm{Res})$
is logically \emph{equivalent} to $\mathsf P\ne\mathsf{NP}$. Theorem~\ref{thm:collapse} adds that every
proof of it must be white-box and nonrelativizing, matching Baker--Gill--Solovay~\cite{BGS75}.
The ``weakest sufficient bridge'' is therefore not weaker than the target for this calculus.

\subsection*{Combined system: resolution plus parity (Builder)}
Tseitin formulas are decided in polynomial time by Gaussian elimination, so any useful
$R$ must at least contain linear-algebra reasoning. $\mathrm{Res}(\oplus)$ (ResLin) does this
and is a fragment of $\mathsf{AC}^0[2]$-Frege. Primary-source check (arXiv:2507.23008v2,
Bhattacharya--Chattopadhyay): exponential lower bounds are known for $\mathrm{Res}(\oplus)$ proofs of
depth $O(N^{2-\varepsilon})$ for lifted Tseitin; the abstract states no superpolynomial size lower
bound is known for unrestricted $\mathrm{Res}(\oplus)$. This identifies a concrete candidate $R$ with an
\textbf{OPEN} lower-bound half, and Theorem~\ref{thm:collapse} still applies to its bridge.

\subsection*{White-box route (Lead)}
\textbf{CONDITIONAL, already in the corpus (LSPV-1/357, /361; Cook 1975~\cite{Cook75}):} if
$\SAT\in\mathsf P$ via an algorithm $D$ whose soundness
$\forall F\,\forall x\,(D(F)=0\to\neg F(x))$ is provable in $\mathrm{PV}$, then EF has polynomial-size
refutations, so $\mathrm{Bridge}(\mathrm{EF})$ holds for such $D$. This route is white-box and therefore
consistent with Corollary~\ref{cor:dichotomy}. Missing arrows: (a) $\SAT\in\mathsf P\Rightarrow$ some
correct decider is PV-provably sound (\textbf{OPEN}); (b) superpolynomial EF lower bounds (\textbf{OPEN},
a central open problem of proof complexity).

\subsection*{Resulting dichotomy (Red Team)}
\begin{center}
\begin{tabular}{@{}lll@{}}
\toprule
$R$ & lower-bound half & bridge half \\
\midrule
Res / hypergraph merge, CP, PC, bounded-depth Frege & KNOWN THEOREM & $\equiv\mathsf P\ne\mathsf{NP}$; white-box only\\
$\mathrm{Res}(\oplus)$ & OPEN (depth-restricted known) & OPEN; white-box only\\
EF & OPEN (major) & CONDITIONAL white-box route via PV\\
\bottomrule
\end{tabular}
\end{center}
Every row needs a nonrelativizing white-box argument somewhere. No row reduces P versus NP to a
strictly easier unconditional statement.

\section{Computational control: resolution width vs.\ Gaussian elimination}
Source \texttt{experiments/claude\_wr/tseitin\_width\_vs\_gauss.py}
(SHA-256 \texttt{2ce678ba\ldots 2cdf43ec}); raw
\texttt{results/raw/claude\_wr/2026-11-13\_1858\_TSEITIN\_WIDTH\_GAUSS\_RAW.json}
(SHA-256 \texttt{22ea37ad\ldots 4b72e}).
Chunking: one graph per Python invocation, total cap 25\,s each. A ``False'' at width $w$ is an
exact completed saturation, so it certifies minimum width $>w$; a timed-out level gives only a lower bound.
Gaussian elimination counts every row XOR.

\begin{center}
\begin{tabular}{@{}lrrrrr@{}}
\toprule
graph & vars & clauses & init.\ width & min.\ Res width & Gauss XORs\\
\midrule
$C_6$ & 6 & 12 & 2 & 2 & 5\\
$K_4$ & 6 & 16 & 3 & 4 & 5\\
prism & 9 & 24 & 3 & 4 & 9\\
$K_{3,3}$ & 9 & 24 & 3 & 5 & 11\\
$K_5$ & 10 & 40 & 4 & $\ge 7$ & 7\\
cube $Q_3$ & 12 & 32 & 3 & $\ge 5$ & 20\\
Petersen & 15 & 40 & 3 & $\ge 5$ & 17\\
grid $3\times3$ & 12 & 32 & 4 & $\ge 5$ & 16\\
\bottomrule
\end{tabular}
\end{center}
Gaussian elimination detects unsatisfiability on all eight. These finite numbers illustrate, but do not
prove, the separation between the hypergraph calculus and a polynomial algorithm. Asymptotic content
comes only from~\cite{BSW01,Urq87}. This is the standard hostile control that already falsifies any
\emph{family-local} (pointwise) bridge for resolution; it does not refute $\mathrm{Bridge}(\mathrm{Res})$,
whose hypothesis is global.

\section*{CORPUS REVIEWED}
\texttt{AGENTS.md}; directive \texttt{research/coordination/2026-11-13\_1249\_TM\_RESOURCE\_AGENT\_REPLACEMENT.md};
the five newest TMR handoffs (TMR-5 BG5 biconnected Tseitin + Gaussian-elimination counteralgorithm);
\texttt{proof1.txt} header (28{,}714 bytes) and size of \texttt{proof2.txt} (1{,}051{,}292 bytes, not
re-read in full this run); Codex 11-11 \emph{SAT polynomial assumption audit} (bridge (B), $B_Q(R)$, local
vs.\ global transfer, ``once $\neg$QBound($R$) is proved, (12) is equivalent to the separation'');
LSPV-1/336, /357, /361 (clocked reflection, PV-to-EF conditional route); repository grep for prior
oracle/black-box bridge statements. Root \texttt{main.tex} was not edited (not a role-owned writer).

\section*{DUPLICATION / CONTRADICTION AUDIT}
Duplicated and adopted, not claimed: Codex 11-11 equivalence of the bridge with the separation once the
lower bound is known; pointwise-bridge falsification by Tseitin (TMR-5/BG5, Codex 11-11); EF/PV conditional
route (LSPV-1/357, /361). New here: the exact equivalence of Theorem~\ref{thm:collapse} via the trivial
oracle $A=\SAT$, and Corollary~\ref{cor:dichotomy}. No contradiction with the corpus was found.

\section*{WHAT CHANGED RELATIVE TO ROOT MAIN}
OPEN status unchanged. One additional PROVED/EXACT barrier constrains every bridge proof.

\section*{ADOPTED / REJECTED / RETIRED PRIOR COMPONENTS}
Adopted: the components listed above. Rejected: any bridge strategy that treats the solver as an
input--output oracle with a step count (Corollary~\ref{cor:dichotomy}). Retired: none.

\section*{STALE OR OVERSTATED CLAIMS FOUND}
None found. Caution: describing $\mathrm{Bridge}(R)$ as ``weaker'' than $\mathsf P\ne\mathsf{NP}$ is
misleading for every $R$ with a known lower bound.

\section*{Barrier transfer and post-quantum ledger}
Logic/bounded arithmetic: Cook 1975 PV$\to$EF is a \emph{directly transferable theorem} (conditional use only).
Relativization/algebrization: directly transferable, and it yields Theorem~\ref{thm:collapse}.
PQ item 5 (black-box worst-case extraction, LWE lesson): \emph{FAILED TRANSFER}. Black-box extraction is
exactly the form excluded by Corollary~\ref{cor:dichotomy}. PQ items 1--3 (Schur squares, quadratic
relations, syzygies): \emph{ANALOGY} only. Polynomial-calculus degree is the natural algebraic analogue,
but Gaussian elimination again shows hidden linear structure is an algorithmic weakness.

\section*{Next targets, one per assignment}
Lead: state the weakest white-box hypothesis (beyond PV-provable soundness) that yields
$\mathrm{Bridge}(\mathrm{Res}(\oplus))$, or show none exists short of EF.
Builder: fix $R=\mathrm{Res}(\oplus)$ with charged linear-form cost; certify polynomial Tseitin refutations.
Hypergraph: find a 3-CNF family candidate-hard for $\mathrm{Res}(\oplus)$ without depth restriction, and
test it against Gaussian elimination and treewidth dynamic programming.
Resource auditor: charge the Cook translation (PV proof $\to$ EF proof) size, including evaluation of $D$.
Red team: try to find a relativizing argument hidden inside any proposed bridge proof.

\begin{thebibliography}{9}
\bibitem{CR79} S.~Cook, R.~Reckhow, The relative efficiency of propositional proof systems, \emph{JSL} 44 (1979) 36--50.
\bibitem{Cook75} S.~Cook, Feasibly constructive proofs and the propositional calculus, \emph{STOC} 1975.
\bibitem{BGS75} T.~Baker, J.~Gill, R.~Solovay, Relativizations of the P=?NP question, \emph{SIAM J.\ Comput.} 4 (1975).
\bibitem{AW09} S.~Aaronson, A.~Wigderson, Algebrization: a new barrier in complexity theory, \emph{ACM TOCT} 1 (2009).
\bibitem{Urq87} A.~Urquhart, Hard examples for resolution, \emph{J.\ ACM} 34 (1987).
\bibitem{BSW01} E.~Ben-Sasson, A.~Wigderson, Short proofs are narrow---resolution made simple, \emph{J.\ ACM} 48 (2001).
\bibitem{BC25} S.~K.~Bhattacharya, A.~Chattopadhyay, Exponential lower bounds on the size of ResLin proofs of nearly quadratic depth, arXiv:2507.23008v2 (2025).
\end{thebibliography}

\noindent P versus NP remains OPEN.
\end{document}
% END_OF_VERIFIED_REPORT::CLAUDE::WR-BRIDGE-0001
